\documentclass[11pt]{article}
\usepackage[a4paper,margin=25mm]{geometry}
\usepackage[T1]{fontenc}
\usepackage{lmodern,amsmath,amssymb,mathtools,graphicx}
\usepackage[hidelinks]{hyperref}
\setlength{\emergencystretch}{2em}
\newcommand{\C}{\mathbb C}
\newcommand{\Tr}{\operatorname{Tr}}
\newcommand{\norm}[1]{\lVert#1\rVert}
\title{The late heat direction through the original observation}
\author{Split-Zero research programme}
\date{21 September 2026}
\begin{document}\maketitle
The supplied smallest-singular-value proof identifies one late heat
scale of the full original operator. We identify its surviving line
as an explicit polynomial class, with a finite error, and carry it
through the unchanged conductor observation. Thus the new full-source
heat calculation has an exact receiver in the original kernel.
The independently attained observed action retains its full coupling.

\section{One specified polynomial class controls the late heat}
Use the complete original polynomial $Q_k$ and attained metric $G_N$
in S1--S45 of the supplied proof \cite{Smallest}. The original
Gamma comparison and reference zero law are
\cite[NG2--3, NG6--16]{NG}; the inherited allowance profile remains
the H4--5 theorem with its stated eventual domain.
All finite statements below precede that asymptotic input.
Let $J_N$ be the full remainder map,
$L_N=H_N^{-1}J_N^*G_N$ its complete minimum section, and
$M[P]=[yP]$ on $E=\C[y]/(Q_k)$. The adjoint is in $G_N$.
Since $k$ is odd and $\delta>0$, $Q_k(0)\ne0$. Define
\[
u=J_N\!\left(\frac{Q_k-Q_k(0)}y\right),\quad
\varphi_N(x)=(L_Nx)(0),\quad
\mathscr Dp=\frac{p-p(0)}y,\quad
B_N=J_N\mathscr D L_N.
\tag{LS1}
\]
The vector $u$ is the same algebraic quotient class at every
cutoff, although its norm and $\varphi_N$ depend on $G_N$.
Multiplying by $y$ gives $Mu=-Q_k(0)[1]$; hence $u\ne0$.
For every $x\in E$, multiplying $B_Nx$ by $M$ gives
$x-\varphi_N(x)[1]$. Therefore
\[
M^{-1}=B_N-\frac{u\varphi_N}{Q_k(0)},\quad
\norm{B_N}\le D_N,\quad
E_N=\frac{\norm u\,\norm{\varphi_N}}{|Q_k(0)|}.
\tag{LS2}
\]
The full explicit $D_N$ is S16. Its proof from the Gamma
bidiagonal inverse and the full original comparison is preserved
in \cite{Smallest}, S11--16; no covariance or source direction
has been deleted.

Assume the finite, directly specified guard $E_N>2D_N$.
The triangle inequality and the kernel of $\varphi_N$ give
\[
\sigma_1(M^{-1})\ge E_N-D_N>D_N\ge\sigma_2(M^{-1}).
\tag{LS3}
\]
Thus the smallest singular value $s_{q,N}$ of $M$ is simple.
Let $P_{\min,N}$ project onto its right singular line, and let
\[
P_{u,N}x=u\frac{u^*G_Nx}{u^*G_Nu},\qquad
\eta_N=\frac{D_N}{E_N-D_N}<1.
\]
Then
\[
\boxed{\norm{P_{\min,N}-P_{u,N}}\le\eta_N,\qquad
\norm{P_{\min,N}-P_{u,N}}_1\le2\eta_N.}
\tag{LS4}
\]
To prove this, choose a unit right singular vector $b$ for the largest
singular value $\sigma$ of $M^{-1}$, and put
$a=M^{-1}b/\sigma$. Applying $I-P_u$ to LS2 kills its rank-one
part and gives $\norm{(I-P_u)a}\le D_N/\sigma\le\eta_N$.
The left singular vector $a$ of $M^{-1}$ is the right singular
vector of $M$ for its smallest value. In the span of $a,u$,
the difference of the two rank-one orthogonal projections has
eigenvalues $\pm\sqrt{1-|\langle a,u/\norm u\rangle|^2}$;
this follows by writing $a=\alpha u/\norm u+\beta w$
in an orthonormal pair. Its operator and trace norms prove LS4.

For any $\tau>0$, the entire spectrum consequently satisfies
\[
\boxed{\left\|e^{-\tau M^\dagger M}
 -e^{-\tau s_{q,N}^2}P_{u,N}\right\|_1
 \le 2\eta_N e^{-\tau s_{q,N}^2}
 +(q-1)e^{-\tau/D_N^2}.}
\tag{LS5}
\]
Indeed the other $q-1$ singular values are at least $D_N^{-1}$
by S5, and their positive spectral projectors contribute the
second term. The first term is exactly LS4 multiplied by the
remaining scalar heat weight. This proof supplies the actual
line; a statement about the smallest singular value alone does
not supply LS5.

\section{The original kernel and its complementary observation}
Let $\Lambda:E\to B$ be the unchanged original observation,
$K=\ker\Lambda$, and
\[
I_K=V^{-1}\pi^T,\quad Z=[M_A,S_kZ_k],\quad \pi Z=0,\quad
P_K=I_K(I_K^*G_NI_K)^{-1}I_K^*G_N .
\tag{LS6}
\]
These are exactly \cite[MR1--10]{MR}. Retain
\[
Q_N=(\Lambda G_N^{-1}\Lambda^*)^{-1},\quad
\mathcal L_N=G_N^{-1}\Lambda^*Q_N,\quad
P_B=\mathcal L_N\Lambda=I-P_K .
\]
Here $\mathcal L_N$ is the observation's minimum section, whereas
$L_N$ in LS1 lifts the full polynomial remainder. Their domains
are $B$ and $E$, respectively. Direct substitution proves
$\Lambda\mathcal L_N=I_B$,
$\mathcal L_N^*G_N\mathcal L_N=Q_N$, and
$\Lambda=\mathcal L_N^{\dagger}$ for these metrics.

The exact kernel fraction of the surviving class is
\[
\begin{split}
\theta_{K,N}
&=\frac{u^*G_N I_K(I_K^*G_NI_K)^{-1}I_K^*G_Nu}{u^*G_Nu}\\
&=1-\frac{(\Lambda u)^*Q_N(\Lambda u)}{u^*G_Nu}\in[0,1].
\end{split}\tag{LS7}
\]
The second equality is the orthogonal decomposition $P_K+P_B=I$.
It is an exact expression on the full original kernel and
the actual observed class $\Lambda u$.

Taking traces in LS5, first against $P_K$ and then against $P_B$,
gives, with its unchanged right side denoted $\epsilon_N(\tau)$,
\[
\begin{gathered}
\left|\Tr(P_K e^{-\tau M^\dagger M})
       -e^{-\tau s_{q,N}^2}\theta_{K,N}\right|
 \le\epsilon_N(\tau),\\
\left|\Tr_B\!\left(\Lambda e^{-\tau M^\dagger M}\mathcal L_N\right)
       -e^{-\tau s_{q,N}^2}(1-\theta_{K,N})\right|
 \le\epsilon_N(\tau).
\end{gathered}\tag{LS8}
\]
The second trace equality uses cyclicity:
$\Tr_B(\Lambda T\mathcal L_N)=\Tr_E(P_BT)$.
The estimate $|\Tr(PT)|\le\norm T_1$ for an orthogonal
projection follows by expanding $T$ into singular vectors and
using $|\langle b,Pa\rangle|\le1$ on each term.

The S4--S5 asymptotics show that $E_N/D_N$ tends to infinity
on the scale $\exp\{q\log(q/k)+O_h(q)\}$; this also follows
from LS2 and $\norm{M^{-1}}=s_{q,N}^{-1}$.
Thus $\eta_N\to0$, and the guard in LS3 holds eventually.
At the common physical time S43, the smallest heat factor tends
to one at $N=q-1,q$ and to zero at $N=2q-1,2q$.
For the exact four-return with signs $+,+,-,-$, LS8 proves
\[
\boxed{\begin{aligned}
\mathcal R\Tr(P_K e^{-\tau_k^*M^\dagger M})
 &=\theta_{K,q-1}+\theta_{K,q}+o_h(1),\\
\mathcal R\Tr_B(\Lambda e^{-\tau_k^*M^\dagger M}\mathcal L_N)
 &=2-\theta_{K,q-1}-\theta_{K,q}+o_h(1).
\end{aligned}}\tag{LS9}
\]
No limit for the individual $\theta$ is assumed. Their sum
cancels exactly between the two original receivers, recovering
the full return $2+o_h(1)$. LS7 is the finite next calculation
for the allocation between these two receivers.

\section{The exact relation to the observed action's own heat}
Let $A_B=\Lambda M\mathcal L_N$. Use the orthogonal decomposition
$E=K\oplus\mathcal L_NB$ with the actual metrics. Write the blocks
of $M$ as
\[
M=\begin{pmatrix}A&D\\ C&A_B\end{pmatrix},\qquad
D=P_KM\mathcal L_N,\quad C=\Lambda M|_K,
\]
so that $H=M^\dagger M$ has blocks
\[
H_{KK}=A^\dagger A+C^\dagger C,\quad
H_{KB}=A^\dagger D+C^\dagger A_B,\quad
H_{BB}=D^\dagger D+A_B^\dagger A_B.
\tag{LS10}
\]
Every term is in the original attained metrics.
For real $z>0$, block elimination gives the complete receiver
\[
\boxed{\Lambda(z+H)^{-1}\mathcal L_N
=\left[z+A_B^\dagger A_B+D^\dagger D
-H_{BK}(z+H_{KK})^{-1}H_{KB}\right]^{-1}.}
\tag{LS11}
\]
The inverse exists because $z+H$ is positive definite; its
Schur complement is positive definite as a complete minimum.
The upper block equation solves the kernel variable as
$-(z+H_{KK})^{-1}H_{KB}$ times the observed variable;
substitution proves LS11.

For completeness, the time-domain connecting equation also retains
these terms. Put $T_B(t)=\Lambda e^{-tH}\mathcal L_N$.
Solving the kernel block with zero initial kernel value gives
\[
\boxed{T_B'(t)=-(A_B^\dagger A_B+D^\dagger D)T_B(t)
+\int_0^tH_{BK}e^{-(t-s)H_{KK}}H_{KB}T_B(s)\,ds,\quad T_B(0)=I.}
\tag{LS12}
\]
Differentiation verifies the equation directly. In finite dimension
its Laplace transform converges for $z>0$ and yields LS11.
Thus LS8--9 return the full original heat to the actual observation;
LS10--12 specify the complete coupling required to compare it with
$e^{-tA_B^\dagger A_B}$. Neither coupling has been set to zero,
and no projected-current phase is inferred from these positive traces.

\section{The surviving class is detected by the actual conductor}
The full kernel in LS6 is contained in the conductor kernel
$V_A=\{[p]:\chi'\mid\mathcal T_Ap\}$ of KF6 \cite{KF}.
Here $\chi$ and $\chi'$ are the complete upper and lower polynomials
of OR1 \cite{OR}; $\chi'$ again denotes the lower polynomial.
Let $\mathcal J$ be the indices of nonzero actual coefficients
$a_j=a_{rs}$, let $J=|\mathcal J|$, and retain the distinct shifts
$b_j=4+(2r-8)\delta+i(2s-8)\gamma$.
Define the actual rational function
\[
\mathcal R_A(z)=\sum_{j\in\mathcal J}\frac{a_j}{z+b_j}
=\frac{P_A(z)}{B_A(z)},\quad
B_A(z)=\prod_{j\in\mathcal J}(z+b_j),\quad
P_A(z)=\sum_{j\in\mathcal J}a_j\prod_{\ell\ne j}(z+b_\ell).
\tag{LS13}
\]
Expansion at infinity gives
$\mathcal R_A(z)=\sum_{n\ge0}(-1)^n\mu_n z^{-n-1}$.
The original symbol order $v$ is at most $J-1$ by the
Vandermonde argument on these $J$ shifts. Consequently
\[
\deg P_A=d:=J-v-1,\qquad
[z^d]P_A=(-1)^v\mu_v\ne0 .
\tag{LS14}
\]
Every residue at $-b_j$ equals the original nonzero $a_j$.
No denominator or zero of this rational function is inserted
by changing the original conductor.

In the original $S$ coordinate the representative of $u$ is
$\widehat u(S)=u((S-c)/i)$. For every lower root $\eta=\beta'_{ab}$,
each $\eta+b_j$ is an upper root, so LS1 gives exactly
\[
(\mathcal T_A\widehat u)(\eta)
=-iQ_k(0)\,\mathcal R_A(\eta-c).
\tag{LS15}
\]
The denominators are nonzero: $\eta+b_j-c$ is a centered
upper root, whose real part is a nonzero odd multiple of $\delta$.
The $q'=(k-7)^2$ arguments $\eta-c$ are distinct. Thus at least
$q'-d$ of these conductor values are nonzero whenever $q'>d$.
In particular $u\notin V_A$, hence $u\notin K$, and
$\theta_{K,N}<1$ at every cutoff on this finite domain.
This is a statement about the original conductor rows, independent
of the unknown metric allocation.

A quantitative version uses the same original rows. Put
$D=\sqrt{\delta^2+\gamma^2}$, $A_1=\sum_j|a_j|$, and impose
$k\ge d+8$. Choose the $d+1$ lower roots
$\eta_j=c'+(2j-k+8)\delta-i(k-8)\gamma$, $0\le j\le d$.
The $d$th finite difference of $P_A$ on this row gives
\[
\sum_{j=0}^d(-1)^{d-j}\binom dj P_A(\eta_j-c)
=(2\delta)^d d!(-1)^v\mu_v.
\]
The sum of the absolute binomial coefficients is $2^d$.
At least one index $j$ therefore satisfies
$|P_A(\eta_j-c)|\ge\delta^d d!|\mu_v|$.
Each denominator factor in LS13 has modulus at most $kD$,
because it is exactly a centered upper root. Equation LS15 gives
\[
\max_{0\le j\le d}|(\mathcal T_A\widehat u)(\eta_j)|
\ge |Q_k(0)|\,\frac{\delta^d d!|\mu_v|}{(kD)^J}.
\tag{LS16}
\]

Let $G_N^{\rm val}=V^{-*}G_NV^{-1}$ be the original attained
metric in the literal upper-root value coordinates, and let
$\kappa_N^{\rm val}$ be its condition number in the Euclidean
metric on those specified coordinates. A conductor row in LS16
is a linear functional $\ell$ on $E$ that annihilates the entire
$K$. Its coefficient vector has Euclidean norm at most $A_1$.
Hence $\norm\ell_{G_N}^2\le A_1^2/\lambda_{\min}(G_N^{\rm val})$.
Also $u(\beta_\alpha)=-iQ_k(0)/(\beta_\alpha-c)$, so
\[
\norm u_{G_N}^2
\le\lambda_{\max}(G_N^{\rm val})q|Q_k(0)|^2/\delta^2.
\]
Since $\ell(u)=\ell(P_Bu)$, the dual norm inequality and LS16
prove the explicit native observation lower bound
\[
\boxed{1-\theta_{K,N}
\ge
\frac{\delta^{2d+2}(d!|\mu_v|)^2}
 {q\,\kappa_N^{\rm val}\,A_1^2(kD)^{2J}}>0.}
\tag{LS17}
\]
The factor $|Q_k(0)|^2$ cancels between the two displayed
physical norms, as the derivation shows. The complete invariant
image is retained: a conductor row annihilates $K$ precisely
because it is a row of its full original observation, not because
the other invariant rows were omitted.
KF33--35 and the original full-form comparison give
$\log\kappa_N^{\rm val}=O_h(q\log(q+2))$.
Thus LS17 is a quantitative positive detection bound at all
four cutoffs. It does not assign a nonzero limiting fraction
or the still uncomputed $kq$ determinant coefficient.

\begin{thebibliography}{9}
\bibitem{Smallest} \emph{The unique small singular direction of the
native arithmetic quotient}, supplied mathematical continuation,
21 September 2026, S1--45. Complete proof included in this edition;
the original correspondence remains private.
\bibitem{NG} Split-Zero research programme,
\emph{Gaussian spectral actions through the original arithmetic
quotient}, NG1--16,
\href{https://github.com/KokunoYumeto/zeta-function-research-reader/blob/5992ad50c0f2a06921f1fcaded7b4e38097c0739/workbenches/splitzero-tandem/continuations/20260920-original-kernel-mixed-determinants/providers/NATIVE_GAUSSIAN_TRANSFER.tex#L101}{public proof, NG6}
and
\href{https://github.com/KokunoYumeto/zeta-function-research-reader/blob/5992ad50c0f2a06921f1fcaded7b4e38097c0739/workbenches/splitzero-tandem/continuations/20260920-original-kernel-mixed-determinants/providers/NATIVE_GAUSSIAN_TRANSFER.tex#L207}{NG10 and its complete proof}.
\bibitem{MR} Split-Zero research programme,
\emph{The original kernel as a mixed determinant}, MR1--23,
\href{https://github.com/KokunoYumeto/zeta-function-research-reader/blob/5992ad50c0f2a06921f1fcaded7b4e38097c0739/workbenches/splitzero-tandem/continuations/20260920-original-kernel-mixed-determinants/ORIGINAL_KERNEL_MIXED_DETERMINANT.tex#L36}{public proof, original frame and metric}.
\bibitem{KF} \emph{The actual invariant kernel: degree filtration,
quantitative minors, and its complete conductor graph}, KF1--55.
Complete proof accompanies this edition in
\href{ACTUAL_KERNEL_FILTRATION_AND_GRAPH.tex}{the included LaTeX source},
particularly KF6--10 and KF33--35.
\bibitem{OR} \emph{The roots of the complete conductor polynomial},
OR1--25. Complete proof accompanies this edition in
\href{COMPLETE_CONDUCTOR_ROOT_GEOMETRY.tex}{the included LaTeX source}.
\end{thebibliography}
\end{document}
